% !TeX root = ../../main.tex
\subsection{Correlated count bounds and the complete normalization ledger}\label{note:zk-count-budget}

\paragraph{Result and boundary.} Bounding every intermediate exponent separately substantially overestimates the memory routers. The exact affine transfer in Proposition~\ref{prop:zk-reused-count-frames} has a second cancellation: total-memory normalization removes both JUMP's bytecode uncertainty and any common shift of its three memory exponents from their remaining correction widths. This note gives uniform interval formulas, smaller repetition schedules and a complete row/label ledger that feeds Corollary~\ref{cor:zk-common-coarse-budget}. It does not yet prove the incoming intervals or row totals for every aggregation witness, and it is not a statistical-transcript theorem.

\subsubsection{Normalize differences, not independent absolute exponents}

Use incoming integer exponents $\zkReuseBcInput{t}$ and $\zkReuseMemInput{t}{c}$, with $t\in\{\mathrm X,\mathrm M,\mathrm S,\mathrm D,\mathrm J\}$. The memory-column counts are $(3,3,1,3,3)$. Define
\[
 \zkNormDiff{t}{c}=\zkReuseMemInput{t}{c}-\zkReuseBcInput{t},
 \qquad
 \zkNormDelta{c}=\zkReuseMemInput{\mathrm J}{c}-\zkReuseMemInput{\mathrm J}{0},
 \quad \zkNormDelta{0}=0.
\]
The required public intervals bound the five bytecode exponents, the ten non-JUMP differences, the single reference difference $\zkNormDiff{\mathrm J}{0}$ and the two other JUMP-column differences. They may contain negative integers. Public constants from existing libraries can be retained as offsets; their sizes do not enlarge interval widths. Compression products must already be public, and no new normalization cycle shares a counted address with the incoming trace.

\paragraph{Unequal centers apply to every stage.} For the centers of Lemma~\ref{lem:zk-unequal-centers}, write
\[
 \zkNormCost{\zkCalibrationCap}=2\zkCalibrationLarge+8\zkCalibrationSmall,
 \qquad
 \zkNormOffset{\zkCalibrationCap}
 =\zkCalibrationLarge(\zkCalibrationLarge-1)
  +4\zkCalibrationSmall(\zkCalibrationSmall-1).
\]
For any deficit from zero through $\zkCalibrationCap$, ten repeated-loop blocks have this fixed row count and add its square offset plus the deficit to the selected bytecode exponent. Reuse one frame per block. Thus an incoming bytecode interval becomes its public upper endpoint plus the square offset at its interval width. The proof of Proposition~\ref{prop:zk-reused-count-frames} uses only equality of each loop's code and operand contributions, so its affine formula remains exact with these unequal centers. Normalize the four non-JUMP bytecode columns before JUMP, including their closing returns.

Write $\zkReuseBcTarget{t}$ for those public normalized bytecode exponents and define
\[
 \zkNormH=\sum_{t\ne\mathrm J,c}\zkNormDiff{t}{c}
              +\sum_{c=0}^{2}\zkNormDelta{c},
 \qquad
 \zkNormS=\zkNormH+3\zkNormDiff{\mathrm J}{0}.
\]
The sum of all thirteen noncompression memory exponents after bytecode normalization is $\zkNormS$ plus the public sum of the corresponding $\zkReuseBcTarget{t}$, one copy per memory column. Summing the supplied intervals gives bounds $\underline{\zkNormH},\overline{\zkNormH}$ and
\[
 \underline{\zkNormS}=\underline{\zkNormH}+3\underline{\zkNormDiff{\mathrm J}{0}},
 \qquad
 \overline{\zkNormS}=\overline{\zkNormH}+3\overline{\zkNormDiff{\mathrm J}{0}},
 \qquad
 \zkNormCap=\left\lfloor\frac{\overline{\zkNormS}-\underline{\zkNormS}}3\right\rfloor.
\]
Install the grouped routers at transfer zero, retaining their public contributions as offsets. Write the remaining deficit $\overline{\zkNormS}-\zkNormS$ as three times its quotient plus its residue in $\{0,1,2\}$. Represent the quotient by the unequal-center construction with cap $\zkNormCap$. Use its ten frames at one common JUMP instruction, so that instruction's total visits are public. The two existing DEREF residue cycles provide the residue in their local-result column. Total memory exponent is now public, with added square offset $3\zkNormOffset{\zkNormCap}$.

\begin{lemma}[JUMP's common component cancels exactly]\label{lem:zk-jump-count-cancellation}
Excluding the grouped routers' fixed contributions, each JUMP memory exponent after this normalization is
\[
 \zkNormPost{\mathrm J}{c}
 =\zkReuseBcTarget{\mathrm J}+\zkNormOffset{\zkNormCap}
  +\zkNormDelta{c}
  +\left\lfloor\frac{\overline{\zkNormS}-\zkNormH}{3}\right\rfloor.
\]
Consequently its correction width depends on the non-JUMP differences and the JUMP-column differences, not on the width of $\zkNormDiff{\mathrm J}{0}$ or of the incoming JUMP bytecode exponent separately.
\end{lemma}
\begin{proof}
Before total-memory normalization this column equals $\zkReuseBcTarget{\mathrm J}+\zkNormDiff{\mathrm J}{0}+\zkNormDelta{c}$. Its new memory labels contribute $\zkNormOffset{\zkNormCap}+\lfloor(\overline{\zkNormS}-\zkNormS)/3\rfloor$. Substitute $\zkNormS=\zkNormH+3\zkNormDiff{\mathrm J}{0}$ and move the integral reference difference through the floor. It cancels exactly. The residue cycles use disjoint JUMP controls, each read once, so add no JUMP memory exponent. Varying the public interval for the reference changes only common offsets, not the resulting interval widths.
\end{proof}

For an explicit tighter interval, put $\zkNormHWithout{c}=\zkNormH-\zkNormDelta{c}$ and sum all the other coordinate bounds. Lower and upper endpoints for the nonconstant part of this JUMP column are
\[
 \left\lfloor\frac{\overline{\zkNormS}+2\underline{\zkNormDelta{c}}-\overline{\zkNormHWithout{c}}}{3}\right\rfloor,
 \qquad
 \left\lfloor\frac{\overline{\zkNormS}+2\overline{\zkNormDelta{c}}-\underline{\zkNormHWithout{c}}}{3}\right\rfloor.
\]
This retains the occurrence of the selected coordinate inside the total rather than charging it twice as independent uncertainty. Non-JUMP intervals are simply those of $\zkNormDiff{t}{c}+\zkReuseBcTarget{t}$, enlarged by at most two at the upper endpoint for DEREF's local-result column.

If $\zkNormHWidth$ and $\zkNormDeltaWidth{c}$ are the widths of these $\zkNormH$ and $\zkNormDelta{c}$ intervals, the displayed JUMP interval has width exactly
\[
 \zkNormColumnWidth{\mathrm J}{c}
 =\zkNormDeltaWidth{c}
   +\left\lfloor\frac{\zkNormHWidth-\zkNormDeltaWidth{c}}3\right\rfloor
   -\left\lfloor\frac{\zkNormDeltaWidth{c}}3\right\rfloor.
\]
Indeed, substituting the summed upper bound for $\overline{\zkNormS}$ makes the two endpoints respectively $\underline{\zkNormDelta{c}}+\overline{\zkNormDiff{\mathrm J}{0}}+\lfloor\zkNormDeltaWidth{c}/3\rfloor$ and $\overline{\zkNormDelta{c}}+\overline{\zkNormDiff{\mathrm J}{0}}+\lfloor(\zkNormHWidth-\zkNormDeltaWidth{c})/3\rfloor$, apart from their common square/code offset. Thus all router widths are computable from interval widths alone, regardless of the public offsets.

\subsubsection{Rows, frames and maximum labels in one ledger}

Let $\zkNormWidth{t}$ be the incoming bytecode interval widths. The selected bytecode rows and total-memory rows are
\[
 \zkNormBcRows{t}=\zkNormCost{\zkNormWidth{t}}\ (t\ne\mathrm J),
 \qquad
 \zkNormBcRows{\mathrm J}=\zkNormCost{\sum_t\zkNormWidth{t}},
 \qquad
 \zkNormMemRows=\zkNormCost{\zkNormCap}.
\]
The sum in the JUMP argument includes the four earlier stages' complementary deficits, not their possibly large public offsets. Choose each integer $\zkNormRouter{t}\ge1$ so that its square bounds the largest postnormalization column width in table $t$, omitting the absorber \texttt{MUL.cnt\_a}. Then the complete noncompression normalizer adds
\[
\begin{aligned}
 \mathrm X:&\quad \zkNormBcRows{\mathrm X}+\zkNormRouter{\mathrm X},\\
 \mathrm M:&\quad \zkNormBcRows{\mathrm M}
  +3(\zkNormRouter{\mathrm X}+\zkNormRouter{\mathrm M}+\zkNormRouter{\mathrm D}+\zkNormRouter{\mathrm J})+\zkNormRouter{\mathrm S},\\
 \mathrm S:&\quad \zkNormBcRows{\mathrm S}+\zkNormRouter{\mathrm S},\\
 \mathrm D:&\quad \zkNormBcRows{\mathrm D}+\zkNormRouter{\mathrm D}+2,\\
 \mathrm J:&\quad \sum_t\zkNormBcRows{t}+\zkNormMemRows
  +\zkNormRouter{\mathrm X}+\zkNormRouter{\mathrm M}+\zkNormRouter{\mathrm S}
  +\zkNormRouter{\mathrm D}+2\zkNormRouter{\mathrm J}+2.
\end{aligned}
\]
Here each grouped router has one selected opcode row, one MUL adapter per routed operand and a closing JUMP. A JUMP-target router therefore adds two JUMPs per traversal. The two residue cycles account for the final twos. All counts are public. After raising every nonabsorbing column to its upper endpoint, the absorber is public by subtracting the other products' integer exponents from the public total. Every compression row and label stays unchanged.

These stages need $50$ bytecode frames, five router frames, ten total-memory frames and four preinstalled residue frames: $69$ frames of at most $128$ cells. The existing templates reserve $175$ code positions, including their unused gaps. For the new labels, take the maximum of the following explicit bounds: the largest traversal count minus one in each bytecode stage; $2\max_t\zkNormRouter{t}-1$ for routers; $\zkNormMemRows-1$ for the common memory-normalization instruction; and one for the residue gadgets. A bytecode stage at cap $\zkCalibrationCap$ has traversal count at most $2\max(\zkCalibrationLarge,\zkCalibrationSmall)$. Zero-length chains contribute no label. All addresses are fresh, so no missing previous count enlarges these maxima.

Add this row vector to the public incoming totals and apply Corollary~\ref{cor:zk-common-coarse-budget}, checking its inequalities and the maximum label $165888$. The two stages together reserve at most $175+308=483$ code positions after the existing $99$ codes, still within log $19$, and conservatively $35+48+3=86$ free slots for these normalizers, routers and final fillers. This excludes the incoming completion, real execution and compression libraries. Their public resource bounds remain necessary; the displayed formulas make that obligation explicit rather than discharging it.

\subsubsection{Unequal-size routers trade spare target rows for fewer MULs}

\begin{lemma}[Rectangular count transfers]\label{lem:zk-rectangular-count-transfer}
A selected memory column with $\zkNormRouter{t}$ reads and its MUL absorber with $\zkNormAdapters{t}{c}$ reads of the same cell can transfer any integer from zero through $\zkNormRouter{t}\zkNormAdapters{t}{c}$ above the selected column's minimal exponent $\binom{\zkNormRouter{t}}2$. The final counter is unchanged, and every label is below $\zkNormRouter{t}+\zkNormAdapters{t}{c}$.
\end{lemma}
\begin{proof}
Use the increasing-sequence construction in Lemma~\ref{lem:zk-square-label-transfer}, now taking the complement inside a chain of length $\zkNormRouter{t}+\zkNormAdapters{t}{c}$. The quotient of the desired transfer by $\zkNormRouter{t}$ is at most $\zkNormAdapters{t}{c}$; equality forces zero remainder. Hence all selected labels remain in this chain, their sum has the required increment, and their complement has exactly $\zkNormAdapters{t}{c}$ labels. Their union is still the complete chain.
\end{proof}

\paragraph{A valid shared-frame implementation.} Put the selected opcode in its own cycle, using a closing JUMP unless it is already a self-loop. In the same frame, give each selected operand a separate cycle of eight MUL instructions that read it as their first input and multiply by zero. The MUL zero-input and output cells are disjoint from the selected operands and have public total accesses. Full traversals supply the quotient of $\zkNormAdapters{t}{c}$ by eight; a separate cycle of at most seven MULs supplies its public remainder. Full and remainder cycles have different return controls, so their destinations coexist in write-once memory. The selected cycle and these absorber cycles need not be traversals of one common path: each is a closed valid component of the state bus. All labels are assigned jointly afterwards. Target operands at offsets below three, zero/output helpers below $38$, target return controls at $64,65,66$, and absorber return controls at $68,\ldots,76$ and $80,\ldots,88$ fit one $128$-cell frame. Every column's shared cell is distinct, so its transfer is independent.

For a public postnormalization width $\zkNormColumnWidth{t}{c}$, choose any public target length $\zkNormRouter{t}\ge1$ and set
\[
 \zkNormAdapters{t}{c}
 =\max\left(1,\left\lceil\frac{\zkNormColumnWidth{t}{c}}{\zkNormRouter{t}}\right\rceil\right),
 \qquad
 \zkNormAdapterTotal=\sum_{t,c}\zkNormAdapters{t}{c},
\]
where the sum runs over the twelve nonabsorbing memory columns. The complete noncompression normalizer now adds
\[
\begin{aligned}
 \mathrm X:&\quad\zkNormBcRows{\mathrm X}+\zkNormRouter{\mathrm X},\\
 \mathrm M:&\quad\zkNormBcRows{\mathrm M}+\zkNormRouter{\mathrm M}+\zkNormAdapterTotal,\\
 \mathrm S:&\quad\zkNormBcRows{\mathrm S}+\zkNormRouter{\mathrm S},\\
 \mathrm D:&\quad\zkNormBcRows{\mathrm D}+\zkNormRouter{\mathrm D}+2,\\
 \mathrm J:&\quad\sum_t\zkNormBcRows{t}+\zkNormMemRows+\sum_t\zkNormRouter{t}
  +\sum_{t,c}\left\lceil\frac{\zkNormAdapters{t}{c}}8\right\rceil+2.
\end{aligned}
\]
This is an alternative to the preceding square-router ledger, not an additional bank. There are still $69$ normalizer frames. Each absorber uses at most seventeen code locations, including both returns; the five target cycles use fourteen including their existing gaps. With bytecode and residue stages, at most $366$ positions suffice. Add $308$ for coarse routers and final fillers, giving $674$ after the old code reservation. For maximum labels replace the square-router term by $\max_{t,c}(\zkNormRouter{t}+\zkNormAdapters{t}{c}-1)$. Absorber-only helper chains and both return chains are shorter. Their fixed contributions do not enlarge the input intervals, and the common-component cancellation still applies.

\paragraph{A large-interval common contract.} The following is a uniform conditional feasibility certificate, not a claim about all real guest witnesses. Suppose the initial noncompression row totals, already including real work, initial filling and earlier libraries, are
\[
 (400000,320000,120000,300000,330000).
\]
Suppose their bytecode widths are bounded by $(10^8,10^8,10^8,10^8,23\cdot10^8)$; each XOR difference width is at most $2\cdot10^8$, each other non-JUMP difference width at most $10^8$, the JUMP reference width at most $25\cdot10^8$, and each of its two nonzero column-difference widths at most $4\cdot10^7$. All interval endpoints must be public, including the offsets from earlier libraries. Extend narrower intervals to these public widths. Set target lengths to
\[
 (24576,16384,16384,32768,49152).
\]
The new normalizer then adds $(45704,118016,37512,53898,451785)$ rows. The shared coarse routers and final fillers fit exactly: their final target-filler lengths are $(4856,12544,293048,96662,20161)$, with $25446$ extra returns from the four non-JUMP fillers. The largest new normalization label is $111441<165888$. Thus all row inequalities and the new-label bound hold for every valid incoming completion satisfying this contract, independently of its interval offsets. By comparison, the previous square-router ledger would require $592559$ MUL rows after coarse routing, exceeding height $19$. The unequal-size construction removes that ledger's bottleneck; it does not certify the incoming trace, memory allocation or statistical privacy.

\begin{lemma}[Batch the target cycles too]\label{lem:zk-batched-count-targets}
In the rectangular construction, each non-JUMP target can use sixteen successive instructions with identical operands, all reading the same selected cells, followed by one return. A separate cycle of at most fifteen target instructions supplies its public remainder. This replaces its $\zkNormRouter{t}$ returns by $\lceil\zkNormRouter{t}/16\rceil$, without changing its transfer capacity, frame count or maximum-label bound.
\end{lemma}
\begin{proof}
Repeat the full cycle the quotient number of times and traverse the remainder cycle once. Use controls $64,65,66$ for the full return and $96,97,98$ for the remainder, disjoint from all absorber controls. All target accesses still share the same cells, so Lemma~\ref{lem:zk-rectangular-count-transfer} assigns exactly the same complete label chains. Each new code visit count and control-chain length is public. The two cycles require at most $17+16=33$ code locations and still fit one $128$-cell frame. The JUMP target remains a self-loop.
\end{proof}

Thus only the JUMP line of the rectangular row ledger changes, to
\[
 \sum_t\zkNormBcRows{t}+\zkNormMemRows+\zkNormRouter{\mathrm J}
 +\sum_{t\ne\mathrm J}\left\lceil\frac{\zkNormRouter{t}}{16}\right\rceil
 +\sum_{t,c}\left\lceil\frac{\zkNormAdapters{t}{c}}8\right\rceil+2.
\]
The whole normalizer now needs at most $486$ code locations, still with $69$ frames. The code cost is public and does not increase the uncertain intervals. In particular it is distinct from batching initial fillers, whose lengths can be private.

\begin{proposition}[A revised common count budget]\label{prop:zk-revised-count-budget}
Suppose the incoming noncompression totals are $(390000,335000,230000,300000,380000)$. Supply public bytecode intervals of widths at most $(150000000,150000000,200000000,150000000,1600000000)$, XOR difference widths at most $250000000$ each, other non-JUMP difference widths at most $200000000$ each, JUMP reference width at most $1600000000$, and JUMP-column difference widths at most $518\cdot65536$ and $867\cdot65536$. Extend narrower intervals to these widths. Use batched target lengths
\[
 (32768,20480,155648,122880,110592).
\]
Assuming the earlier compression normalization, disjoint reservations and coarse-backdrop hypotheses, and incoming labels at most $165888$, the all-column normalizer, shared coarse routers and final fillers fit all five tables. The new normalizer's maximum label is $156932$.
\end{proposition}
\begin{proof}
The exact interval and row formulas give added counts $(58510,115350,185276,148624,443540)$. Substituting in Corollary~\ref{cor:zk-common-coarse-budget} leaves target-filler lengths $(2050,210,35284,1936,1382)$ and exactly $2470$ extra non-JUMP filler returns. Every residual is nonnegative and the new labels obey the cap. The same targets without batching require $311040$ further returns and exceed JUMP capacity.
\end{proof}

The initial totals are realized by the private-length recipe in \noteref{zk-count-prefill}, under its stated real-row caps. That note reserves the necessary code and frames and derives both JUMP width bounds from the frame contract. The four other bytecode ranges, the remaining memory differences and the real indirect-label cap still need uniform certification. This revised contract therefore addresses the previously exhibited incoming MUL mismatch without asserting full candidate feasibility or extending the statistical view.

\subsubsection{Private-length prefilling need not create quadratic differences}

\begin{lemma}[Linear difference cost for the initial fillers]\label{lem:zk-prefill-count-differences}
Use the sixteen-instruction filler of Lemma~\ref{lem:zk-batched-local-filler} for a possibly private length $16\zkNormFillQuotient+\zkNormFillRemainder$, $0\le\zkNormFillRemainder<16$. In every memory column of its selected opcode it changes $\zkNormDiff{t}{c}$ by exactly $-\zkNormFillQuotient\zkNormFillRemainder$. If the selected opcode is not JUMP, its closing returns change all three JUMP differences $\zkNormDiff{\mathrm J}{c}$ by $-\zkNormFillQuotient\mathbf 1_{\zkNormFillRemainder>0}$. Neither kind changes a JUMP-column difference $\zkNormDelta{c}$.
\end{lemma}
\begin{proof}
Each target memory column contributes $16\binom{\zkNormFillQuotient}{2}$, because the suffix uses a second fresh frame. Its bytecode column contributes this same value plus $\zkNormFillQuotient\zkNormFillRemainder$, since the suffix revisits exactly $\zkNormFillRemainder$ of the same instruction addresses. Their difference has absolute value at most $15\zkNormFillQuotient$. For a non-JUMP filler, the single return instruction contributes $\binom{\zkNormFillQuotient+\mathbf 1_{\zkNormFillRemainder>0}}{2}$, whereas each return memory column contributes $\binom{\zkNormFillQuotient}{2}$. Their difference is the stated value, equal in all three columns. A JUMP-only filler is covered by the target calculation. Additivity follows from fresh addresses between different filler blocks.
\end{proof}

Thus completing private real-row counts to public initial totals can be budgeted through small difference intervals even when absolute code and memory sums are quadratic. One must still bound the private filler lengths, their bytecode intervals and all their closing returns before fixing the initial JUMP total. This lemma neither makes those intermediate sums public nor hides their transcript disclosures.

\subsubsection{Account for BLAKE2s-priority relabeling explicitly}

At each real memory cell, stably move its BLAKE2s accesses before its other accesses, preserving the old label order inside both groups. A noncompression access gains exactly the number of BLAKE2s accesses whose old label was larger than its own. All gains are nonnegative. Every cell still uses the same complete label chain, so the total memory exponent and all final counters are unchanged; the total gain outside BLAKE2s equals its exponent loss. Bytecode counts do not change. Consequently, if a memory role can share its cell with at most $\zkPriorityCap{t}{c}$ BLAKE2s accesses, this relabeling shifts that column's exponent, and its memory-minus-code difference, by an integer in $[0,\zkPriorityCap{t}{c}\zkPriorityRows{t}]$, where $\zkPriorityRows{t}$ bounds its real rows.

Under the same fresh-frame, owned-local-operand and single-instruction-visit contract as the static BLAKE2s load bound, the public function bodies bound $\zkPriorityCap{t}{c}$ for every local role. DEREF's indirect role instead uses the largest BLAKE2s cell load. The actual optional local-range recursive wrapper passes the following checked upper bounds, in its native memory-column order:
\[
\begin{array}{c|rrr}
 t&\multicolumn{3}{c}{\zkPriorityCap{t}{c}}\\\hline
 \mathrm X&2&16&16\\
 \mathrm M&2&2&1\\
 \mathrm S&213&\cdot&\cdot\\
 \mathrm D&0&213&182\\
 \mathrm J&16&1&2
\end{array}
\]
In particular the two JUMP-column differences shift within $[-16\zkPriorityRows{\mathrm J},\zkPriorityRows{\mathrm J}]$ and $[-16\zkPriorityRows{\mathrm J},2\zkPriorityRows{\mathrm J}]$, rather than intervals based on the global $256$ cap for every role. The research allocator example checks these public incidences before execution; its native two-child execution still passes. This is a conditional static bound for the relabeling increment, not a bound on the original differences or a proof that every hinted execution obeys the ownership contract. Apply it before the disjoint initial fillers and include all earlier count libraries in the incoming intervals.

\subsubsection{Local-role priority bounds the original JUMP differences}

The preceding increment bound retains the original chronological memory counts as a premise. A different canonical assignment removes that dependence for local roles. At every real memory address, keep the BLAKE2s accesses first, then place \texttt{JUMP.cnt\_c}, \texttt{JUMP.cnt\_d} and \texttt{JUMP.cnt\_f}, in that order. Put all remaining local roles next, in any fixed public role order, and indirect \texttt{DEREF.cnt\_target} last. Within a role retain the previous order. Apply this only to the real execution; disjoint padding libraries keep their own assignments. In particular the BLAKE2s labels agree exactly with the earlier stable compression-priority assignment.

\begin{lemma}[Local counts do not inherit foreign indirect loads]\label{lem:zk-local-role-priority}
Suppose each real call owns a disjoint frame, every local operand stays in that frame, and its executed local occurrences form a subset of those in its owning public function, without repetition. For a noncompression local role $(t,c)$, let $\zkLocalCountCap{t}{c}$ be the maximum, over public functions and operand occurrences in that role, of the number of local occurrences at the same relative cell in that or an earlier priority class, minus one. Include all BLAKE2s occurrences, but exclude indirect DEREF targets. The new label of each such access is at most $\zkLocalCountCap{t}{c}$, regardless of the number or destinations of indirect reads. Consequently
\[
 0\le\zkReuseMemInput{t}{c}
 \le\zkLocalCountCap{t}{c}\zkPriorityRows{t}.
\]
If $\zkLocalReads{a}$ and $\zkIndirectReads{a}$ are the real local and indirect read counts at cell $a$, the remaining indirect exponent is exactly
\[
 \zkReuseMemInput{\mathrm D}{1}
 =\sum_a\left(\zkLocalReads{a}\zkIndirectReads{a}
                    +\binom{\zkIndirectReads{a}}2\right).
\]
All memory values, final counters, code labels and the total memory exponent are unchanged.
\end{lemma}
\begin{proof}
Assign the complete chain by the prescribed priority. A local label counts only earlier local occurrences. Ownership and single visits inject these occurrences into those counted in the public function body; omitting a private branch cannot increase their number. The indirect labels at a cell are exactly $\zkLocalReads{a},\ldots,\zkLocalReads{a}+\zkIndirectReads{a}-1$, whose sum is the displayed expression. Each address still has the same complete chain, so legality and the conserved quantities follow. This includes aliased operand roles. The statement does not bound indirect multiplicity through frame size.
\end{proof}

\paragraph{Two incoming widths now follow from the actual public bytecode.} The optional local-range ordinary-child wrapper statically certifies
\[
 (\zkLocalCountCap{\mathrm J}{0},\zkLocalCountCap{\mathrm J}{1},\zkLocalCountCap{\mathrm J}{2})
 \le(517,1,350).
\]
Therefore its original, real-execution differences satisfy the public intervals
\[
 \zkNormDelta{1}\in[-517\zkPriorityRows{\mathrm J},\zkPriorityRows{\mathrm J}],
 \qquad
 \zkNormDelta{2}\in[-517\zkPriorityRows{\mathrm J},350\zkPriorityRows{\mathrm J}].
\]
Their widths are at most $518\zkPriorityRows{\mathrm J}$ and $867\zkPriorityRows{\mathrm J}$. A real JUMP budget of $40000$ gives widths $20720000$ and $34680000$, both below the earlier $4\cdot10^7$ contract; both remain below that cap through $46136$ real JUMPs. These are program-wide conditional bounds, not intervals centered on measured exponents. Lemma~\ref{lem:zk-prefill-count-differences} preserves them under the initial sixteen-instruction fillers. Earlier disjoint libraries with public JUMP-difference increments contribute public offsets only; any library with private increments must still be charged. In particular $330000$ in the incoming row ledger includes padding and must not be substituted for the real-row premise.

For the other local differences, a public bytecode interval immediately gives
\[
 \zkNormDiff{t}{c}\in\left[-\overline{\zkReuseBcInput{t}},
 \zkLocalCountCap{t}{c}\zkPriorityRows{t}-\underline{\zkReuseBcInput{t}}\right]
\]
before the disjoint libraries and fillers. These bounds can be loose: constant pooling creates large static incidences, and the indirect column still depends on the displayed collision sum. This lemma resolves the two original JUMP-difference intervals under a concrete real-row cap, not the five bytecode intervals, the full incoming capacity contract or joint statistical privacy. Ownership and uniform real-row bounds remain proof obligations for the intended recursive program class.

\paragraph{Certificates.} The budget audit checks eight complete ISA libraries with privately varying foreign reads into JUMP controls, aliased local/indirect accesses and compression cycles. Priority routing preserves all complete chains, BLAKE2s priority, bytecode exponents and the total memory exponent, and it verifies the exact indirect formula. The native allocator audit additionally reconstructs this label assignment for the real wrapped guest, including resolved indirect addresses, checks every real final counter against the native memory image, and checks its total against both the original trace and the complete-chain identity. This is a label-legality and count-profile check, not a native padded proof. The static JUMP caps are checked before running the witness and therefore do not depend on its private trace.

\paragraph{Exact certificates and the remaining premise.} \auditref{zk_count_budget_audit.py} checks all small signed interval/residue cases, exact filler remainders, the row/code/frame/label ledger and eight complete cycle libraries with independent memory reuse and instruction multiplicities. A MUL reading a JUMP control cell also gives genuinely unequal, privately varying JUMP-column sums. A separate private choice changes only the incoming JUMP bytecode exponent; the three postnormalization JUMP exponents and the residue remain exactly unchanged, as the cancellation predicts. All final column products and memory/code images agree. Its \texttt{--coarse} mode composes four of these libraries with the shared coarse routers and checks every native height-seven count node. All ISA constraints and full read chains pass, with compression unchanged.

As an algebraic stress case, widen JUMP's bytecode interval to $[0,2147450880]$ and its reference-difference interval to $[-2147450880,70]$, retaining the small fixture's other intervals. The router sizes remain $(3,4,3,3,6)$; the normalizer adds $(13,71,13,15,190301)$ rows and its maximum new label is $95111$. This is an exact conditional ledger, not an aggregation resource theorem: the fixture's small non-JUMP and JUMP-column difference bounds are not being attributed to the native guest. Next certify those invariant intervals after BLAKE2s-priority label assignment and initial public-total filling, include every existing library and apply the common residual inequalities. Full joint GKR, other reductions and authentication remain open independently of this capacity progress.

The \texttt{--rectangular} mode checks the same private families with unequal target and absorber lengths, including both directions of imbalance and batched remainder cycles; combining it with \texttt{--coarse} checks the resulting native frontier. The audit also certifies the large-interval contract above, including the failed square-router comparison, and checks stable BLAKE2s-priority relabeling on eight complete cycle libraries with aliased operands and varied prior label orders. These are legality and exact-budget certificates. They neither assert that the native guest satisfies the large-interval contract nor extend the proved statistical view.
